\slajdid{09_3-s055}
\begin{frame}[label=09_3-s055]{Pravilna veza prolaznih tranzistora}
% element: 09_3-s055:e01
% element: 09_3-s055:e02
% element: 09_3-s055:e03
% element: 09_3-s055:e04
\begin{columns}[c,onlytextwidth]\column{.47\textwidth}
\begin{center}
{\RazaviSetup
\begin{circuitikz}[font=\sitneoznake,line width=.7pt]
\ctikzset{transistors/scale=.85,bipoles/length=.55cm}

\foreach \x/\y/\i/\v in {0/0/1/B,1.5/1/2/C,3/2/3/D}{
\node[rz nmos,nobulk,rotate=90](t\i)at(\x,\y){};
\draw[wire](t\i.D)--(\x-.6,\y)node[ocirc]{}node[left]{$\v$};\node[above]at(\x,\y+.2){\i};}
\draw[wire](t1.G)--++(0,-.3)node[ocirc]{}node[below]{$A$};
\draw[wire](t1.S)--(1.5,0)|-(t2.G);\node[above]at(.7,0){$V_1$};
\draw[wire](t2.S)--(3,1)|-(t3.G);\node[above]at(2.2,1){$V_2$};
\draw[wire](t3.S)--(3.8,2)node[right]{$Y$};\draw(3.6,2)to[C,l=$C_L$](3.6,1)node[ground]{};\node[junction]at(3.6,2){};
\draw[DEcrvena](-.6,-.5)--(3.8,2.7);\draw[DEcrvena](-.6,2.7)--(3.8,-.5);
\end{circuitikz}}

\end{center}
\column{.49\textwidth}
\begin{center}
{\RazaviSetup
\begin{circuitikz}[font=\sitneoznake,line width=.7pt]
\ctikzset{transistors/scale=.85,bipoles/length=.55cm}

\foreach \x/\i/\v in {0/3/C,1.3/2/B,2.6/1/A}{
\node[rz nmos,nobulk,rotate=90](t\i)at(\x,1){};
\draw[wire](t\i.G)--++(0,-.3)node[ocirc]{}node[below]{$\v$};\node[above]at(\x,1.2){\i};}
\draw[wire](t3.D)--(-.6,1)node[ocirc]{}node[left]{$D$};
\draw[wire](t3.S)--(t2.D);\draw[wire](t2.S)--(t1.D);
\draw[wire](t1.S)--(3.1,1)--(3.1,.35)--(4.1,.35)node[ocirc]{}node[above]{$Y$};
\draw[wire](3.1,.35)--(3.1,0);\node at(3.1,-.4){$\vdots$};\node[junction]at(3.1,.35){};
\draw(3.9,.35)to[C,l=$C_L$](3.9,-.65)node[ground]{};\node[junction]at(3.9,.35){};
\end{circuitikz}}

\end{center}
\end{columns}\[V_1=V_{DD}-V_{Tn1},\qquad V_2=V_{DD}-V_{Tn1}-V_{Tn2}\]\[Y=V_{DD}-V_{Tn1}-V_{Tn2}-V_{Tn3}\]\noindent Za $V_A=V_B=V_C=V_D=V_{DD}$ pravilna veza daje $V_Y=V_{DD}-V_{Tn}$.
\par Ostatak mreže realizujemo selektorski da izbegnemo neodređeno stanje i visoku impedansu.
\end{frame}
